\frontmatter
\thispagestyle{empty}

\includegraphics[height=2.2cm]{dept_of_methodology}
\hfill
\includegraphics[height=2.2cm]{lse}

\vspace*{3cm}
\begin{center}
{\LARGE \textbf{MY451}}\\[2ex]{\LARGE \textbf{Introduction to Quantitative Analysis}}
\\[7ex]
%{\Large Michaelmas Term 2005}\\[15ex]
{\Large 2015-16}\\[15ex]
{\Large Jouni Kuha}\\[2ex]{\Large Department of Methodology}\\[1ex]{\Large London School of Economics and
Political Science}
\end{center}

\newpage
\thispagestyle{empty}

$\, $
\newpage
\setcounter{page}{1}


%\setcounter{chapter}{-1}
\chapter{Course information}


{\Large \textbf{MY451: Introduction to Quantitative Analysis}}


\subsubsection{Course Description}
This course is intended for those with little or no past training in
quantitative methods.  The course is an intensive introduction to some
of the principles and methods of statistical analysis in social
research.
Topics covered in MY451 include descriptive statistics, basic ideas of
inference and estimation, contingency tables and an introduction to
linear regression models.
For those with some quantitative training the slightly more
advanced course MY452 (Applied Regression Analysis) would be more
appropriate, followed by other Department of Methodology and Department
of Statistics courses on quantitative methods, such as MY454 (Applied Statistical Computing), MY455 (Multivariate Analysis and
Measurement), MY456 (Survey Methodology), MY457 (Causal Inference for
Observational and Experimental Studies),
MY459 (Quantitative Text Analysis),
ST416 (Multilevel Modelling), and ST442 (Longitudinal Data Analysis).

Please note that up to 2010-11 the course was called MI451, so older
materials will bear that code. For example, you will find past
examination papers on the Library website under MI451.

\subsubsection{Course Objectives}
This course aims to impart a level of familiarity suitable for a
moderately critical understanding of the statistical material in the
journals commonly used by students in their work and the ability to use
some elementary techniques.


\subsubsection{Staff}
\begin{itemize}
\item
\textbf{Teachers responsible}: Dominik Hangartner in MT;
Jonathan Jackson in LT
\item
\textbf{Other Lecturers}:  Benjamin Wilson (MT) and Matthew Mulford (LT)
\item
\textbf{Course Administrator}: Esther Heyhoe, email
\texttt{e.heyhoe@lse.ac.uk}
\item
\textbf{Class Teachers}:
See the MY451 Moodle site for the list of class teachers.
\end{itemize}

\newpage
%\subsubsection{Location}
\subsubsection{Teaching}
\begin{itemize}
\item
\textbf{Lectures}:
Ten two-hour lectures in weeks 1 - 5 and 7 - 11:
\begin{itemize}
\item
Michaelmas Term (MY451M): Thu 2.00pm - 4.00pm in the Peacock Theatre
\item
Lent Term (MY451L): Wed 2.00pm - 4.00pm in the Hong Kong Theatre
(CLM.G.02)
\end{itemize}
\item
\textbf{Computer classes}:
Each student will attend a one-hour class each week,
\textbf{starting in Week 2}. Please see on-line timetables for the times
and places of the class groups.
The allocation of students to classes is done through
LSE for You or the Moodle system (see below), depending on your status. This will be
exaplained in the first lecture and on the MY451 Moodle page.
Please do not change the classes allocated to you without our consent.
\end{itemize}

\subsubsection{Course Materials}
\begin{itemize}
\item
\textbf{Coursepack}:
This coursepack is the main course text.
Printed copies of this coursepack will be available for purchase at the
beginning of each term, and
instructions on how to obtain
a hard copy will be given at the first lecture and on the MY451 Moodle
page.
Digital copies of this coursepack are available
for download for free from the Moodle page.
\item
\textbf{Lecture slides}: Copies of most of the slides displayed during
the lectures can be downloaded from the MY451 Moodle page.
\item
\textbf{Recommended course texts}:
\begin{itemize}
\item
Alan Agresti and Christine Franklin (2013). \emph{Statistics: The Art and
Science of Learning from Data} (Third Ed.). Pearson.
\item
Alan Agresti and Barbara Finlay
(2013). \emph{Statistical Methods for the Social Sciences} (Fourth Ed.).
Pearson
\end{itemize}
Earlier editions are also suitable.
While neither of these books is absolutely required, we
recommend that you purchase one if you would like to have additional
explanation, examples and exercises to supplement the coursepack. Please
note that during the examination you will be allowed to use the
coursepack and any other notes or books that you
find useful. Of these two, Agresti and Finlay is a bit more advanced. It
is also the recommended additional course
text for MY452 (which also has a coursepack similar to this one),
so you may want to purchase it if you are planning to also take
that course.
\item
\textbf{Other text books}: There are hundreds of introductory textbooks
on statistics and quantitative methods, many of them covering almost
identical material. If you have one which you would like to use, and
which looks as if it covers the same material at about the same
level as this course, then it is probably suitable as
additional reading. If in doubt, please ask one of the lecturers.
\begin{itemize}
\item
There are also many books which focus on the SPSS statistical software
package used in the computer classes. We do not consider them necessary
for this course, or for learning statistics or SPSS. If, however, you
feel otherwise, some such books are mentioned on page
\pageref{p_spssbooks}.
\end{itemize}
\end{itemize}

%All registered LSE students will be provided with ONE course pack free
%of charge.

%Supplementary Reading
%Robert Perry and John D. Robertson (2002), Comparative Analysis of Nations: Quantitative Approaches. Westview Press.
%Champney,L Introduction to Quantitative Political Science, Harper Collins College Publishers, 1995.  This book gives multiple examples of how statistical techniques are commonly used in political science.  Also,   Pennings, Paul, Hans Keman \& Kleinnijenhuis Doing Research in Political Science: An Introduction to comparative methods and statistics by.


\subsubsection{Homework}
\vspace*{-1.5ex}
There will be homework assignments each week. These are your opportunity
to practice interpreting the results of the types of analyses that are introduced in
the lectures and then carried out in the computer classes. The homework questions
are also broadly similar in style to typical examination questions for
MY451. The questions are based on results of the computer class
exercises or of further similar analyses. For the homework, you can
use output that you produced yourself in the class or afterwards.
Alternatively, all the computer outputs that are
required to answer the homework questions can be downloaded from the
Moodle page.

Homework should be submitted by Monday after its distribution in class, at the latest
by 5pm in MT and by 10am in LT. Please see the MY451 Moodle page for
instructions on how to submit your homeworks.
Homework submitted in time will be corrected and returned at the next
class. Please make sure to clearly indicate your name, and the time,
class number and class teacher of your computer class on all homework
assignments. For the homework assignments, students are
encouraged to work together.

\subsubsection{MY451 on Moodle}
\label{sss_moodle}
The course materials are all available on Moodle. Go to
\texttt{http://moodle.lse.ac.uk/} and login using your \emph{username}
and \emph{password} (the same as for your LSE e-mail). Then in the
\emph{select courses} dialogue box type in MY451, and in \emph{search
results} click on MY451. The site contains the structure of the course
week by week, the readings, weekly computer class assignments and the
associated data sets, coursepack and other materials, as well as a
section on news and announcements.

\subsubsection{Advisory Sessions}
There will be a time each week during the Michaelmas and Lent terms during which a
member of the teaching team will be available at the Department of
Methodology in Columbia House to answer
questions about the course. Information on the times of the advisory
hours will be given at the first lecture. These sessions are not
intended to be private tutorials and you will be expected to have made a
genuine attempt to work through the relevant sections of the coursepack
prior to coming to an advisory session. Also, questions addressing
material from that week's lecture will be given priority.  If you are
having difficulty with a topic, come and see us at that time. If you
wait until the end of term, when the advisory sessions can be very busy,
we cannot guarantee that your questions from earlier sessions will be
covered.  \textbf{There will be no advisory sessions after the end of
Lent Term.}

\subsubsection{Notes on studying for the course}
To learn the material from this course you must do the work every week
since it is cumulative; if you miss a week or two (or sleep through it!)
there is a chance that you will find yourself hopelessly lost. So this
is definitely not a ``pick and choose'' course!  Also bear in mind that
most people cannot learn quantitative techniques passively by just
turning up to the lectures and reading the occasional chapter in a
textbook.  To learn statistics you have to do it; there are no
shortcuts. Thus in addition to a two-hour weekly lecture there will be
one-hour computer classes (in which you do some data analysis and
interpretation using SPSS - instructions will be provided) and there
will be weekly homework (which will be corrected but not graded by your
class teacher). Doing the assignments in the computer classes and the
homework demonstrate whether you have understood and can apply what was
covered in the lectures.  If you are having any trouble this will reveal
what the problem is.  Thus the course is designed to have multiple,
reinforcing ways of helping you get to grips with this material.


\subsubsection{Examinations/assessment}

\vspace*{-2ex}
There will be a \textbf{two-hour examination in the Summer Term 2016}.
This is an \textbf{open-book examination}, to which you are allowed to
bring in any written material, including the coursepack and any text
book with your notes. You are required to bring your own calculators for
the examination. Examination papers from previous years are available
for revision purposes at the LSE library web site. Students should
understand that past examinations should only be used as rough guides to
the types of questions that are likely to appear on the examination.

For many of you, MY451 is only one part of a package of methods training
with a course code such as MC4M1 or MY4M1. If this is the case, your
result for MY451 will contribute only part of the final mark for the
package, with a weight determined by the regulations.

\vspace*{-1ex}
\subsubsection{Computing}

\vspace*{-2ex}
Students must know their Username and Password in time for the first
class.  This information can be obtained from IT Help Desk (Library, 1st
floor).  The course uses Microsoft Windows-based software.
If you are not familiar with the program, you might want to
attend introductory
courses in Windows during the first two weeks of the term.
The statistical package being used is SPSS, which will be
introduced in the first computing class of the course.

\vspace*{-1ex}
\subsubsection{Software availability}

\vspace*{-2ex}
A personal copy of the program SPSS for Windows can be bought by course
participants (maximum 1 copy per person) for the price of
%approximately
\pounds10
from IT services, \textbf{not} the Department of Methodology.

\vspace*{-1ex}
\subsubsection{Feedback}

\vspace*{-2ex}
We would welcome any comments you have on the course.  If there are any
problems that we can deal with, we will attempt to do so as quickly as
possible.  Speak to any member of the course team, or to your
departmental supervisor if you feel that would be easier for you.
Also please let us know if you find any errors or omissions in the
coursepack, so that we can correct them for next year.

%Substantial parts of this coursepack were rewritten and rearranged for
%2012-13. Many new typographical errors and other flaws will no doubt
%have been introduced in the process. Please do tell us about any that
%you discover.

\vspace*{-1ex}
\subsubsection{Acknowledgements}

\vspace*{-2ex}
The computer classes were extensively revised for 2012-13. The hard work
of designing the new exercises and writing the revised instructions was
done by Sally Stares.

This coursepack bears many traces of previous materials and all of their
authors, Colm O'Muircheartaigh, Colin Mills, Matt Mulford, Fiona Steele,
Paul Mitchell, and Sally Stares. Many thanks to Farimah Daftary, Sue
Howard, Jon Jackson, Paul Mitchell, Indraneel Sircar, and many students
of previous years for comments and suggestions which are incorporated in
the current revision.


\newpage
\subsubsection{Course Programme}

\begin{tabular}{lll}
\multicolumn{2}{l}{\textbf{Week 1}} & \\
& Lecture  & Course overview and organisation. Introduction to basic
concepts\\
& Class & \textbf{No class}\\
& Coursepack & Chapter \ref{c_intro}  \\[3ex]
\multicolumn{2}{l}{\textbf{Week 2}} & \\
& Lecture & Descriptive statistics for categorical variables  \\
& Class & Introduction to SPSS. Descriptive statistics for categorical variables\\
& Coursepack &
Sections
\ref{s_descr1_intro}--\ref{s_descr1_2cat} and
\ref{s_descr1_presentation}\\[3ex]
\multicolumn{2}{l}{\textbf{Week 3}} & \\
& Lecture \& Class & Descriptive statistics for continuous variables \\
& Coursepack & Sections
\ref{s_descr1_1cont}--\ref{s_descr1_2cont}
\\[3ex]
\multicolumn{2}{l}{\textbf{Week 4}} & \\
& Lecture  \& Class & Analysis of two-way contingency tables \\
& Coursepack & Chapters \ref{c_samples} and \ref{c_tables}
\\[3ex]
\multicolumn{2}{l}{\textbf{Week 5}} & \\
& Lecture  \& Class& Inference for means in two populations \\
& Coursepack & Chapters \ref{c_contd} and \ref{c_means} \\[3ex]
\multicolumn{2}{l}{\textbf{Week 6}} & \\
&Reading Week & \textbf{No lecture, no class}. Revision quiz  available on Moodle. \\[3ex]
\multicolumn{2}{l}{\textbf{Week 7}} & \\
& Lecture  \& Class & Inference for proportions in one and two populations\\
& Coursepack & Chapter \ref{c_probs} \\[3ex]
\multicolumn{2}{l}{\textbf{Week 8}} & \\
& Lecture  \& Class & Correlation and simple linear regression as descriptive methods \\
& Coursepack & Sections \ref{s_regression_intro}--\ref{ss_regression_simple_est} \\[3ex]
\multicolumn{2}{l}{\textbf{Week 9}} & \\
& Lectures  & \\
& \hspace*{1em}Hour 1  &
Inference for the simple linear regression model \\
& \hspace*{1em}Hour 2  &
Three-way contingency tables \\
& Class & Simple linear regression \\
& Coursepack & Section \ref{ss_regression_simple_inf}
(Hour 1); Section \ref{s_regression_causality} and
Chapter \ref{c_3waytables} (Hour 2) \\[3ex]
\multicolumn{2}{l}{\textbf{Week 10}} & \\
& Lecture \& Class & Multiple linear regression \\
& Coursepack & Sections \ref{s_regression_multiple}--\ref{s_regression_rest}
\\[3ex]
\multicolumn{2}{l}{\textbf{Week 11}} & \\
& Lecture & Review and further statistical methods\\
& Class & Discussion and additional exercise on multiple linear regression \\
& Coursepack & Chapter \ref{c_more}
\end{tabular}

\newpage
\subsubsection{FAQ: Frequently Asked Questions}

\textbf{Why do we use SPSS? I've heard that SAS/STATA/MINITAB/R/LIMDEP is
much better.} At this level it does not matter which program you use since
we are learning standard procedures that are common to all programs. In
favour of SPSS is that it has an easy to learn interface and
that it is widely available at sites both within and outside the UK.

\textbf{Can I get a copy of the SPSS software to use on my home computer?}
Yes,
for a small fee from IT Services helpdesk, not from the Department of Methodology.

\textbf{I'm taking MY451 because I want to learn how to use SPSS but we don't
seem to learn very much about the program. Why is that?} MY451 is not a
course about learning to use SPSS. We use the program merely to
facilitate data analysis and interpretation.
Some options for learning  more
about SPSS will be mentioned in the first lecture.

\textbf{I'm taking MY451 to help me analyse data for my dissertation. Can I
discuss my data and my specific problems with the lecturers?} Yes, but
not during the course.
Staff of the Department of Methodology
will be happy to talk to you about
problems specific to your dissertation during the weekly sessions of the
Methodology Surgery (see the website of the department for more information).

\textbf{Does the coursepack contain everything I need
to know for the exam?} Yes. However, you will
stand by far the best chance in the exam if you also
attend the lectures, where the lecturers emphasise and explain
the key parts of the material.

% A
%very large proportion of it, but the course pack is not a text-book. Its
%purpose is to save you the trouble of struggling to write down detailed
%notes in the lecture and to enable you to give more attention to what is
%said.

\textbf{The lecturer introduced some material
that was not in the coursepack.
Do I need to know that material?}
This is almost certainly an illusion. The lectures will not introduce
any genuinely new material not included in the course pack. However,
sometimes the lecturer may of course use different words or a
different example to further explain some topic.
Copies of the most relevant notes
displayed at the lectures will be posted in the MY451 Moodle site.
All of the material required for the exam is contained in the
coursepack, with the posted lecture notes as additional clarification.

%The course-pack is a set of notes. It
%is inevitable that it won't contain everything that is said in the
%lecture. It would be boring in the extreme if all the lecturer did was
%simply go through the course-pack word for word. Strictly speaking the
%examinable course consists of the lectures, the coursepack and the
%readings from the textbooks.

%Can I pass the exam without looking at Healey or an equivalent textbook?
%Probably not.

\textbf{Do I need to bring my coursepack to lectures and classes?}
Yes, because some of the material displayed during the lectures will be
graphs or tables straight from the coursepack.
However, an even more important way of using the coursepack is to read
the relevant parts of it \emph{before} each lecture. In addition, you
will find the instructions for the class exercises at the end of the
coursepack.

%I have lost my coursepack: can I get a new one? Each student is entitled
%to one course pack gratis. If you lose your pack you will have to
%download a copy from the MI public folders on MS Outlook or the Mi451
%WebCT page and print it out yourself.

\textbf{I can't make it to the advisory hour. Why aren't there more of them?}
Advisory hours are offered in addition to the normal support
arrangements for School courses and rely on the goodwill and enthusiasm
of the course team. Given that most team members also offer their time
to support advisory hours for other courses there is a limit to how much
time they can volunteer. If you genuinely are unable to attend an
advisory hour, but desperately need advice, lecturers will always see
you in their personal office hours. If these are inconvenient you may be
able to set up an appointment at another time by special arrangement.

\newpage
\textbf{There was a long queue at advisory and I didn't get much/any time to
discuss my particular problem. Isn't there a better way to arrange
things?} Advisory sessions are meant for quick feedback and clarification
of problems on a first-come first served basis. They are not meant to be
personal tutorials, seminars or impromptu lectures. They are also not
meant to be a substitute for attending the lectures; reading the
course-pack; doing the homework; thinking. Priority
is always given to problems related to the previous week's lecture
material. If several people have the same problem it will usually be
possible to talk to everyone together, but there is a finite limit to
how many people can be crammed into an office at the same time!

\textbf{I want to bring my homework to the advisory session because I
couldn't understand the class teacher's comments. Is that OK?} Yes, but
it is unlikely that the person that wrote the comments on your homework
will be taking the advisory hour. Usually it will be much better to
discuss homework comments directly with your class teacher in class.

%\textbf{Can I submit my homework by email?} No, not normally. It takes
%time and effort to print out emailed submissions and there is no
%guarantee that what is printed out on our printers will be formatted in
%the way you intend. However, your class teacher may allow this if he/she
%has a different preference.
%An exception may also be made
%if you are officially registered as a part-time student and
%therefore are only able to attend the School on particular days.

\textbf{Can I work together on the homework with my friends?} Yes, we positively
encourage you to discuss the homework assignments with your colleagues.
If you do this, please submit \textbf{one} copy to be marked rather
than multiple copies of the same answers.

\textbf{If I get less than 50\%
in the exam what happens now?} Candidates sit
MY451 under different examination regulations depending on which
department they are in and what degree they are registered for.  For
some degrees a pass in MY451 is required, for others the mark from MY451
is just one component of an overall course mark (implying a ``pass'' mark
in MY451 is not required). The regulations for your degree are not set
by Methodology, but by your home department. To find out whether a pass in
MY451 is required, consult your degree regulations
or ask the tutor responsible for your program in your home
department. Candidates who fail at the first attempt and whose degree
regulations require a pass are, as per the usual School examination
rules, entitled, at the next available opportunity, to ONE more attempt.
This will be in the Summer Term of the next Session.

\textbf{Why don't we get our examination scripts returned after they
have been marked?} Written work is of two sorts, formative and
summative. The former, for example the homework exercises, is meant to
give you feedback on how you are doing and will, where appropriate, be
returned with comments. The latter, for example the examination, is
meant to evaluate your performance, not explain to you where you have
gone wrong, or for that matter where you have done exceptionally well.

\textbf{If I don't agree with the mark, how do I appeal?} The cultural norm in
the UK is that marks are not arrived at by a process of teacher-student
negotiation. You can make a formal appeal through the School's appeal
process (see the appropriate section of the website for
details). NB: appeals cannot be made on grounds of academic substance,
only on grounds of procedural irregularities. In other words an appeal
will not be allowed if the only grounds you have for appealing is that
you/your friend/your personal advisor/your spiritual guru think your
script was worth more marks than the examiners did.

\textbf{The class teachers are always very busy during class helping
people with problems and I can't get enough personal attention. Why
can't there be more class teachers in the classroom?} For MY451 there
will usually be two class teachers per classroom i.e. double the normal
complement. Even so we realise that you may not be able to get attention
at the exact moment you want it. Please be patient and if help is not
immediately available you can always try asking the person who is
sitting next to you!

\textbf{I'm not registered at the LSE but at another University of London
college. Can I attend this course?} Normally yes, but you will have to
complete an intercollegiate enrolment form.

\textbf{I would like to audit the course without taking the exam. Is that OK?}
Yes, you are welcome to attend the lectures providing you are an
LSE/University of London student and there is room for you.

\textbf{MY451 is not challenging enough for me. Is there a more
difficult course?} Yes, MY452 and numerous other courses offered by the
Department of Methodology and the Statistics department.

\tableofcontents
\mainmatter
